\documentclass[aspectratio=169]{beamer}
\usetheme{Madrid}
\usecolortheme{beaver}
\usepackage[utf8]{inputenc}
\usepackage[T1]{fontenc}
\usepackage{tikz}
\usepackage{amssymb}
\usepackage{graphicx}
\setbeamertemplate{navigation symbols}{}
\setbeamercolor{title}{fg=white,bg=cyan!60!black}
\setbeamercolor{frametitle}{fg=white,bg=cyan!50!black}
\setbeamercolor{block title}{bg=cyan!15,fg=black}
\setbeamercolor{block body}{bg=cyan!5,fg=black}

\title{Story Studio 5 — Choose a Path}
\subtitle{if/else — a story that branches on the reader's choice}
\author{ByteForge --- Grade 7 --- Game Maker --- Week Overview}
\date{}

\begin{document}

\begin{frame}
    \titlepage
\end{frame}

\begin{frame}{This Week's Plan}
    \begin{itemize}
        \item Lesson 1: Two Endings
        \item Lesson 2: Three Paths
        \item Lesson 3: Branch on the Click
        \item[$\bigstar$] BOSS --- Your Story, Your Ending (\texttt{lessonBoss.tex})
    \end{itemize}
\end{frame}

\begin{frame}{Real-World Example}
    \Large
    A "choose your own adventure" book, a Netflix interactive episode (Bandersnatch), or a game's dialogue choice — every branch a reader/player can take is an if/else somewhere in the code. Minecraft villagers branch too: trade with them enough and their dialogue/trades change — a real if/else on a "reputation" variable.
\end{frame}


\begin{frame}{Let Them Discover (1/2)}
    \Large
    Think of a choice-based story or game you know. Sketch its first three branches on paper as boxes and arrows (this is secretly a flowchart) before we build one in code.
\end{frame}


\begin{frame}{Let Them Discover (2/2)}
    \Large
    In Minecraft, trading with a villager enough times unlocks better trades. What do you think the game is checking each time — and what "branch" fires once you cross that threshold?
\end{frame}

\begin{frame}{Files in This Folder}
    \small
    \begin{itemize}
        \item \texttt{lesson1.tex} --- Two Endings
        \item \texttt{lesson2.tex} --- Three Paths
        \item \texttt{lesson3.tex} --- Branch on the Click
        \item \texttt{lessonBoss.tex} --- BOSS: Your Story, Your Ending
    \end{itemize}
    \vspace{0.5em}
    \normalsize Each lesson file is its own short deck --- open the one for
    the specific day you're teaching, not this whole overview.
\end{frame}

\end{document}
